\chapter{Fundamentos Geométricos}

\begin{quotation}
	\itshape
	``Comienza con la métrica de Fisher--Rao'' es la instrucción técnica; ``comienza reconociendo que tu única brújula fiable es la curvatura de tu propia incertidumbre'' es la contraseña filosófica.
\end{quotation}

\section{El Espacio de las Probabilidades}

El punto de partida filosófico de la geometría de la información es la decisión normativa de tratar a la incertidumbre como un objeto geométrico primario. En lugar de considerar las distribuciones de probabilidad simplemente como funciones analíticas o herramientas para calcular esperanzas, las consideramos como \emph{puntos} en un espacio abstracto.

Imaginemos una familia paramétrica de distribuciones de probabilidad $S = \{p_\theta : \theta \in \Theta\}$, donde cada punto en el conjunto $\Theta \subseteq \mathbb{R}^n$ etiqueta una distribución única. Nuestro objetivo es dotar a este conjunto $S$ de una estructura geométrica que refleje las propiedades intrínsecas de la inferencia estadística, y no simplemente la topología euclidiana de los parámetros $\theta$.

La pregunta fundamental es: ¿Qué significa que dos distribuciones estén ``cerca''? La respuesta no la dicta la naturaleza, sino nuestra postura epistémica: la geometría debe ser invariante bajo reparametrizaciones y transformaciones suficientes.

\section{Variedades Estadísticas}

Formalmente, una \textbf{variedad estadística} es una variedad diferenciable $S$ cuyos puntos son distribuciones de probabilidad. Consideramos un espacio muestral $\mathcal{X}$ y una familia de funciones de densidad de probabilidad $p(x; \theta)$ con respecto a una medida base (usualmente Lebesgue o conteo).

\[ S = \{ p(x; \theta) \mid \theta = (\theta^1, \dots, \theta^n) \in \Theta \} \]

Asumimos condiciones de regularidad estándar (suavidad con respecto a $\theta$, soporte común, etc.) que permiten tratar a $\Theta$ como un sistema de coordenadas local para $S$.

\section{La Métrica de Fisher}

Para medir distancias en esta variedad, introducimos la \textbf{Matriz de Información de Fisher}. Esta matriz actúa como el tensor métrico Riemanniano en nuestra variedad estadística.

Dada una función de log-verosimilitud $\ell(\theta; x) = \log p(x; \theta)$, la métrica de Fisher $g_{ij}(\theta)$ se define como la covarianza del score:

\[ g_{ij}(\theta) = \mathrm{E}_\theta \left[ \frac{\partial \ell}{\partial \theta^i} \frac{\partial \ell}{\partial \theta^j} \right] = \int_{\mathcal{X}} p(x; \theta) \frac{\partial \ell}{\partial \theta^i} \frac{\partial \ell}{\partial \theta^j} \, dx \]

Bajo condiciones de regularidad, esto es equivalente a la esperanza del Hessiano negativo de la log-verosimilitud (la curvatura de la sorpresa):

\[ g_{ij}(\theta) = -\mathrm{E}_\theta \left[ \frac{\partial^2 \ell}{\partial \theta^i \partial \theta^j} \right] \]

\begin{definition}[Información de Fisher]
	En dimensión $n = 1$, la métrica de Fisher se reduce a la \textbf{información de Fisher}:
	\[
	I(\theta) = \mathrm{E}_\theta\!\left[\left(\frac{\partial \log p_\theta(X)}{\partial \theta}\right)^2\right] = -\mathrm{E}_\theta\!\left[\frac{\partial^2 \log p_\theta(X)}{\partial \theta^2}\right].
	\]
	En el caso general, la matriz $[I_{ij}] = [g_{ij}]$ actúa como una métrica Riemanniana en el espacio de parámetros.
\end{definition}

\subsection{Ejemplos de cálculo de la métrica de Fisher}

\paragraph{Bernoulli.} Para $X \sim \mathrm{Bern}(p)$ con $p \in (0,1)$:
\[
\log p(x;p) = x \log p + (1-x)\log(1-p), \qquad
\frac{\partial \log p}{\partial p} = \frac{x}{p} - \frac{1-x}{1-p}.
\]
\[
\frac{\partial^2 \log p}{\partial p^2} = -\frac{x}{p^2} - \frac{1-x}{(1-p)^2} \;\;\Longrightarrow\;\;
I(p) = -\mathrm{E}\!\left[\frac{\partial^2 \log p}{\partial p^2}\right] = \frac{1}{p} + \frac{1}{1-p} = \frac{1}{p(1-p)}.
\]
Cuando $p$ se acerca a $0$ o $1$, la información diverge: las distribuciones casi degeneradas son extremadamente sensibles a cambios en $p$.

\paragraph{Poisson.} Para $X \sim \mathrm{Poisson}(\lambda)$ con $\lambda > 0$:
\[
\log p(x;\lambda) = -\lambda + x \log \lambda - \log(x!), \qquad
I(\lambda) = \frac{1}{\lambda}.
\]
A mayor tasa $\lambda$, menor información por observación: las distribuciones con alta media son ``más difíciles'' de distinguir.

\paragraph{Exponencial.} Para $X \sim \mathrm{Exp}(\lambda)$ con $\lambda > 0$:
\[
\log p(x;\lambda) = \log \lambda - \lambda x, \qquad
I(\lambda) = \frac{1}{\lambda^2}.
\]

\paragraph{Normal.} Para $X \sim \mathcal{N}(\mu, \sigma^2)$ con $\sigma^2$ fijo:
\[
I(\mu) = \frac{1}{\sigma^2}.
\]
Si ambos $\mu$ y $\sigma$ son desconocidos, la matriz de información es:
\[
I(\mu, \sigma) = \begin{pmatrix} 1/\sigma^2 & 0 \\ 0 & 2/\sigma^2 \end{pmatrix}.
\]

\begin{center}
	\begin{tikzpicture}
		\begin{axis}[
			width=10cm, height=6cm,
			xlabel={$p$},
			ylabel={$I(p)$},
			domain=0.05:0.95,
			samples=100,
			axis lines=left,
			ymax=20,
			grid=both,
			grid style={gray!30},
			tick style={black},
			title={Información de Fisher: Bernoulli},
		]
			\addplot[Accent, thick] {1/(x*(1-x))};
			\node[AccentDark] at (axis cs:0.5,5) {$I(p) = \frac{1}{p(1-p)}$};
			\addplot[AccentDark, dashed] coordinates {(0.5,0) (0.5,20)};
		\end{axis}
	\end{tikzpicture}
\end{center}

\section{Interpretación Geométrica}

La métrica de Fisher $g$ define un producto interno en el espacio tangente $T_\theta S$ de la variedad en cada punto $\theta$. Si consideramos dos desplazamientos infinitesimales $d\theta$ y $\delta\theta$, su producto interno es:

\[ \langle d\theta, \delta\theta \rangle_g = \sum_{i,j} g_{ij}(\theta) d\theta^i \delta\theta^j \]

Esto nos permite medir la longitud de curvas en la variedad estadística. La distancia entre dos distribuciones $p_{\theta_1}$ y $p_{\theta_2}$ es la longitud de la geodésica (el camino más corto) que las conecta bajo esta métrica.

\begin{center}
	\begin{tikzpicture}[scale=0.85]
		\draw[thick, Accent, fill=AccentLight!12, rounded corners=15pt]
			(-0.5,0.2) .. controls (1.5,-0.4) and (4.5,-0.3) .. (6.5,0.8)
			.. controls (7.2,2) and (5.5,3.5) .. (3.5,3.2)
			.. controls (1.5,3) and (0,2.8) .. (-0.3,1.5)
			.. controls (-0.5,0.8) and (-0.6,0.3) .. (-0.5,0.2);
		\node[AccentDark, font=\small] at (3,1.2) {$\mathcal{M}$ (variedad)};
		\fill[AccentDark] (2,2.6) circle (3pt);
		\node[AccentDark, above=3pt, font=\small] at (2,2.6) {$p_{\theta_1}$};
		\draw[thick, Accent!30] (0.6,2.3) -- (3.2,2.9);
		\node[font=\tiny, Accent!60, right] at (3.2,2.9) {$T_{\theta_1}\mathcal{M}$};
		\draw[->, thick, red!70!black] (2,2.6) -- (2.8,2.8) node[above, font=\scriptsize] {$d\theta$};
		\draw[->, thick, blue!70!black] (2,2.6) -- (1.5,2.8) node[left, font=\scriptsize] {$\delta\theta$};
		\fill[AccentDark] (5,1.5) circle (3pt);
		\node[AccentDark, below=3pt, font=\small] at (5,1.5) {$p_{\theta_2}$};
		\draw[->, AccentDark, thick] (2.2,2.5) .. controls (3.5,2.8) and (4.5,2.2) .. (4.8,1.7);
		\node[AccentDark, font=\small] at (4,2.7) {geodésica};
	\end{tikzpicture}
\end{center}

La distancia, la \textbf{distancia de Fisher--Rao}, cuantifica qué tan distinguibles son dos distribuciones basándose en muestras observadas. Si la información de Fisher es alta, una pequeña perturbación en los parámetros resulta en un gran cambio en la distribución (gran distancia estadística).

\begin{example}[Distancia de Fisher--Rao para $\mu$ en Normal]
	Para $X \sim \mathcal{N}(\mu, \sigma^2)$ con $\sigma^2$ fijo, la geodésica entre $\mu_1$ y $\mu_2$ es una recta, y la distancia es:
	\[
	d(\mu_1, \mu_2) = \int_{\mu_1}^{\mu_2} \sqrt{I(\mu)} \, d\mu = \frac{|\mu_1 - \mu_2|}{\sigma}.
	\]
	La normalización por $\sigma$ captura la intuición: si $\sigma$ es grande (mucho ruido), dos medias cercanas son ``difíciles'' de distinguir.
\end{example}

\section{El caso bivariante: Normal $(\mu, \sigma)$}

Cuando consideramos ambos parámetros $\mu$ y $\sigma$ como variables, la geometría se vuelve notablemente rica. La matriz de Fisher es:

\[
I(\mu, \sigma) = \begin{pmatrix} 1/\sigma^2 & 0 \\ 0 & 2/\sigma^2 \end{pmatrix}.
\]

El elemento off-diagonal es cero: $\mu$ y $\sigma$ son \emph{ortogonales} bajo la métrica de Fisher. Esto significa que, localmente, la información sobre la media es independiente de la información sobre la varianza.

\begin{proposition}[Geometría hiperbólica del semiplano]
	El espacio de parámetros $(\mu, \sigma)$ con $\sigma > 0$, dotado de la métrica de Fisher, es isométrico al \textbf{semiplano superior de Poincaré} $\mathbb{H}^2 = \{(\mu, \sigma) : \sigma > 0\}$ con métrica $ds^2 = \frac{d\mu^2 + 2\, d\sigma^2}{\sigma^2}$. Las geodésicas son semicírculos ortogonales al eje $\sigma = 0$ y verticales.
\end{proposition}

\begin{center}
	\begin{tikzpicture}[scale=0.7]
		\draw[->] (-4,0) -- (4,0) node[right] {$\mu$};
		\draw[->] (0,0) -- (0,4.5) node[above] {$\sigma$};
		\draw[thick, Accent!30, domain=-1.4:1.4, samples=100] plot ({\x}, {sqrt(1.5^2 - \x^2)}) node[right, font=\scriptsize] {geodésica};
		\draw[thick, Accent!50, dashed] (2,0.5) -- (2,4) node[above, font=\scriptsize] {vert.};
		\fill[AccentDark] (-1.2,1.5) circle (2pt) node[above left, font=\scriptsize] {$p_1$};
		\fill[AccentDark] (1.8,1.2) circle (2pt) node[below right, font=\scriptsize] {$p_2$};
		\node[font=\scriptsize, gray] at (3.5,0.5) {$\sigma=0$ (degenerada)};
	\end{tikzpicture}
\end{center}

Las distribuciones con $\sigma$ pequeño (poca varianza) están ``lejos'' de las distribuciones con $\sigma$ grande, aún si sus medias son idénticas. Esta es la geometría hiperbólica subyacente: el eje $\sigma = 0$ es ``el infinito'' de la variedad.

\section{Invariancia y Unicidad}

Un resultado fundamental debido a Chentsov \cite{chentsov1982statistical} establece que la métrica de Fisher es la \emph{única} métrica Riemanniana (salvo factor escalar) en una variedad estadística que es invariante bajo \textbf{estadísticos suficientes}.

\begin{theorem}[Chentsov, 1972]
	Si $g$ es una métrica Riemanniana en una familia de distribuciones que es invariante bajo toda transformación estocástica suficiente (maps de Markov que preservan información), entonces $g$ es proporcional a la métrica de Fisher.
\end{theorem}

\textbf{Idea de la demostración.} Un \emph{map de Markov} $T : \mathcal{X} \to \mathcal{Y}$ transforma datos sin perder información: si $T(X)$ es suficiente para $\theta$, la geometría no debería cambiar. Chentsov muestra que existe un \emph{único} tensor simétrico positivo definido que es invariante bajo todos los maps de Markov suficientes. Dicho tensor es, necesariamente, la métrica de Fisher (hasta una constante).

Esta propiedad de invariancia confirma que la métrica de Fisher captura la estructura intrínseca de la información, independiente de la representación de los datos o de la parametrización elegida.

\section{Conexión con la cota de Cramér--Rao}

La métrica de Fisher tiene una consecuencia práctica directa: establece un límite inferior para la varianza de cualquier estimador insesgado.

\begin{theorem}[Cramér--Rao]
	Para cualquier estimador insesgado $\hat\theta$ de $\theta$:
	\[
	\operatorname{Var}(\hat\theta) \ge \frac{1}{n \, I(\theta)}.
	\]
\end{theorem}

La inversa de la métrica de Fisher es, literalmente, la menor varianza alcanzable. Un modelo con información de Fisher alta (curvatura alta) permite estimación más precisa; uno con información baja, menos. La geometría decide los límites, no la ingeniería del estimador.

\begin{example}[Cota de Cramér--Rao para Bernoulli]
	Para $n$ muestras de $\mathrm{Bern}(p)$, la cota es:
	\[
	\operatorname{Var}(\hat p) \ge \frac{p(1-p)}{n}.
	\]
	El estimador $\hat p = \bar X$ (proporción muestral) alcanza esta cota asintóticamente: es \emph{eficiente}.
\end{example}

\section{Ejercicios}

\subsection*{Conceptos Básicos}
\begin{enumerate}
	\item[$\star$] (V/F) Una variedad estadística es un conjunto de parámetros sin estructura adicional.
	\item[$\star$] (V/F) La métrica de Fisher depende del sistema de coordenadas elegido para los datos $x$.
	\item[$\star$] Defina explícitamente el conjunto $S$ para la familia de distribuciones Bernoulli. ¿Cuál es el espacio de parámetros $\Theta$?
	\item[$\star$] Escriba la función de log-verosimilitud para una sola observación $x$ de una Bernoulli($\theta$).
	\item[$\star$] Explique con palabras qué representa la información de Fisher: ¿qué mide?
\end{enumerate}

\subsection*{Cálculo de Información de Fisher}
\begin{enumerate}[resume]
	\item[$\star$] Calcule la primera derivada (score) de la log-verosimilitud de Bernoulli con respecto a $\theta$.
	\item[$\star$] Calcule la segunda derivada de la log-verosimilitud de Bernoulli.
	\item[$\star$] Calcule la Información de Fisher $I(\theta)$ tomando la esperanza del negativo de la segunda derivada.
	\item[$\star$] Calcule $I(\lambda)$ para la distribución de Poisson $P(x;\lambda) = e^{-\lambda}\lambda^x/x!$.
	\item[$\star$] Calcule $I(\lambda)$ para la distribución Exponencial $p(x;\lambda) = \lambda e^{-\lambda x}$.
	\item[$\star$] Calcule $I(\mu)$ para la Normal $\mathcal{N}(\mu, \sigma^2)$ asumiendo $\sigma^2$ conocido y fijo.
	\item[$\star$] Calcule $I(\sigma)$ para la Normal $\mathcal{N}(0, \sigma^2)$ asumiendo $\mu=0$ conocido.
	\item[$\star\star$] Calcule la matriz de Información de Fisher $2\times2$ para la Normal $\mathcal{N}(\mu, \sigma^2)$ con ambos parámetros desconocidos.
	\item[$\star\star$] Calcule $I(\theta)$ para la distribución de Pareto $p(x;\theta) = \theta x^{-\theta-1}$ para $x \ge 1$.
	\item[$\star\star$] Calcule $I(p)$ para la distribución Geométrica $P(X=k) = (1-p)^{k-1}p$.
	\item[$\star\star$] Calcule $I(\alpha)$ para la distribución $\mathrm{Gamma}(\alpha, \beta)$ con $\beta$ conocido.
	\item[$\star\star\star$] Escriba la expresión para $I(\alpha)$ de una distribución Beta con $\beta$ fijo (en términos de funciones poligamma).
\end{enumerate}

\subsection*{Métrica en 2 parámetros}
\begin{enumerate}[resume]
	\item[$\star\star$] Para $\mathcal{N}(\mu, \sigma^2)$ con ambos parámetros desconocidos, muestre que los parámetros $\mu$ y $\sigma$ son ortogonales bajo la métrica de Fisher.
	\item[$\star\star$] Calcule la distancia de Fisher--Rao infinitesimal entre $\mathcal{N}(0, \sigma_1)$ y $\mathcal{N}(0, \sigma_2)$ (media fija en 0).
	\item[$\star\star\star$] Verifique que la métrica $ds^2 = \frac{d\mu^2 + 2\,d\sigma^2}{\sigma^2}$ corresponde al semiplano de Poincaré con un factor de escala. ¿Es $\mathbb{H}^2$ completo o solo isométrico a una porción?
\end{enumerate}

\subsection*{Cambio de Variables e Invariancia}
\begin{enumerate}[resume]
	\item[$\star\star$] Si $\eta = f(\theta)$ es una reparametrización, ¿cómo se relaciona $I(\eta)$ con $I(\theta)$? (Regla de transformación tensorial 1D).
	\item[$\star\star\star$] Aplique lo anterior a la Bernoulli usando el parámetro logit $\eta = \log(\theta/(1-\theta))$. Calcule $I(\eta)$.
	\item[$\star\star\star$] Verifique que para Poisson con reparametrización $\lambda = e^{\eta}$, la información de Fisher en $\eta$ es igual a la de $\lambda$ multiplicada por $e^{2\eta}$. ¿Qué significa esto geométricamente?
\end{enumerate}

\subsection*{Distancias y Geometría}
\begin{enumerate}[resume]
	\item[$\star\star$] Escriba la fórmula para la distancia de Fisher--Rao entre dos distribuciones parametrizadas por $\theta_1$ y $\theta_2$ en una variedad 1D.
	\item[$\star\star$] Calcule la distancia de Fisher--Rao entre $\mathcal{N}(0, 1)$ y $\mathcal{N}(1, 1)$ (varianza fija).
	\item[$\star\star\star$] Calcule la distancia de Fisher--Rao entre $\mathcal{N}(0, 1)$ y $\mathcal{N}(0, 2)$ (media fija).
	\item[$\star\star\star\star$] (Reto) ¿Cuál es la distancia entre dos distribuciones categóricas (multinomiales) $p$ y $q$? (Pista: geometría esférica).
\end{enumerate}

\subsection*{Propiedades Teóricas}
\begin{enumerate}[resume]
	\item[$\star$] Enuncie la Cota de Cramér--Rao en términos de la Información de Fisher.
	\item[$\star$] ¿Cuál es la varianza mínima asintótica para un estimador insesgado de $p$ en una Bernoulli?
	\item[$\star$] ¿El estimador de máxima verosimilitud alcanza esta cota asintóticamente?
	\item[$\star$] ¿Quién probó que la métrica de Fisher es la única invariante bajo estadísticos suficientes?
	\item[$\star$] ¿Qué divergencia induce localmente la métrica de Fisher?
	\item[$\star\star$] Escriba la expansión de Taylor de segundo orden de la divergencia KL $D_{KL}(p_\theta \| p_{\theta+\delta})$ en términos de $I(\theta)$.
	\item[$\star\star$] ¿Qué representa la curvatura de la función de log-verosimilitud en el máximo?
	\item[$\star\star$] ¿Qué representa la varianza del score?
	\item[$\star\star\star$] ¿Por qué es importante la raíz cuadrada del determinante de la métrica de Fisher en la inferencia Bayesiana?
	\item[$\star\star\star$] Muestre que la cota de Cramér--Rao se alcanza exactamente para el estimador de $\lambda$ en $\mathrm{Poisson}(\lambda)$ con $n$ muestras.
	\item[$\star\star\star$] Verifique que el determinante $\det I(\mu, \sigma)$ para la Normal bivariada es $2/\sigma^4$. Interprete: ¿qué significa que $\det I \to \infty$ cuando $\sigma \to 0$?
\end{enumerate}

\subsection*{Ejercicios avanzados}
\begin{enumerate}[resume]
	\item[$\star$] \textbf{Invariancia tensorial.} (V/F) ¿La métrica de Fisher es invariante bajo reparametrización del parámetro $\theta$? \textit{Pregunta a tu LLM:} «En Bernoulli, si cambio a $\eta=\log(p/(1-p))$, ¿qué le pasa a $I(p)$?»
	\item[$\star\star$] \textbf{Fisher Multinomial.} Para $X\sim\mathrm{Categorical}(\pi_1,\pi_2,\pi_3)$ con $\sum\pi_i=1$, calcule la matriz de Fisher $I(\pi)\in\mathbb{R}^{3\times 3}$ y muestre que $\det I=0$ (la familia vive en el simplex $2$-dimensional). \textit{Pista:} el score es $\partial_{\pi_k}\log p = \mathbb{1}_{X=k}/\pi_k - 1$; use la restricción $\sum\pi_k=1$ como ligadura.
	\item[$\star\star\star$] \textbf{Invariancia tensorial.} Demuestre que bajo $\eta=\log(p/(1-p))$ en Bernoulli se cumple $g^{\eta}_{\eta\eta}=g^p_{pp}(dp/d\eta)^2$, y verifique numéricamente con $p=0.3$ y $p=0.7$. \textit{Verifica con:} $I(\eta)\stackrel{?}{=}(1/(p(1-p)))\cdot(\theta(1-\theta))^2 \cdot (4 e^\eta/(1+e^\eta)^2)$ en $\eta=0$.
	\item[$\star\star\star$] \textbf{Pareto + Cramér-Rao.} Para $X\sim\mathrm{Pareto}(\theta)$ con $p_\theta(x)=\theta x^{-\theta-1}$ ($x\ge 1$), calcule $I(\theta)$ y demuestre que $\hat\theta = n/\sum\log X_i$ alcanza la cota de Cramér-Rao asintóticamente. \textit{Verifica con:} simulación con $n=10^5$, $\theta=2$.
	\item[$\star\star\star\star$] \textbf{Mini-proy: Geodésicas en $\mathbb{H}^2$.} Implemente las tres $\alpha$-geodésicas ($\alpha=-1,0,1$) entre $\mathcal{N}(0,1)$ y $\mathcal{N}(3,2^2)$ sobre $(\mu,\sigma)$. \textit{Hint:} para $\alpha=1$ (exponencial), la geodésica es recta en $\eta=(\mu/\sigma^2,-1/(2\sigma^2))^T$. \textit{Discuta:} ¿cuál cruza $\sigma=0$?
\end{enumerate}

\subsection*{Ejercicios avanzados}
\begin{enumerate}[resume]
	\item[$\star$] (V/F) Toda $f$-divergencia satisface la desigualdad triangular. \textit{Pregunta a tu LLM:} «¿Por qué Hellinger sí pero KL no es métrica?»
	\item[$\star\star$] Para $p=(0.7,0.2,0.1)$ y $q=(0.4,0.4,0.2)$, calcule $D_{\mathrm{KL}}(p\|q)$ y Jensen--Shannon $J(p,q)$. \textit{Pista:} $J(p,q)=\tfrac{1}{2} D_{\mathrm{KL}}(p\|m)+\tfrac{1}{2} D_{\mathrm{KL}}(q\|m)$ con $m=(p+q)/2$.
	\item[$\star\star$] Calcule 5 $f$-divergencias ($f(t)=t\log t$, $-\log t$, $(t-1)^2$, $(1-\sqrt{t})^2$, $|t-1|$) entre dos Bernoulli y grafique la asimetría $D_f(p\|q)-D_f(q\|p)$ para 1000 pares $(p,q)$. \textit{Verifica con:} la única ``métrica verdadera'' (Hellinger) tiene asimetría 0.
	\item[$\star\star\star$] Para $f(t)=(t-1)^2$, demuestre que $D_f(p\|q)=\sum(p-q)^2/q$ (link con $\chi^2$ clásico). \textit{Hint:} sustituya y use $\sum p=\sum q=1$.
	\item[$\star\star\star\star$] \textbf{Mini-proy: KL vs Wasserstein sobre MNIST.} Entrene dos clasificadores: uno con pérdida KL y otro con Wasserstein-2 sobre pares de dígitos $8\times 8$ perturbados. Reporte accuracy y robustez adversaria ante 5 niveles de ruido gaussiano $\sigma\in[0,0.3]$. \textit{Discuta:} ¿en qué la métrica W es más robusta?
\end{enumerate}

\subsection*{Ejercicios avanzados}
\begin{enumerate}[resume]
	\item[$\star\star$] \textbf{Gamma = exponencial.} Descomponga $\mathrm{Gamma}(\alpha{=}3,\beta{=}2)$ en la forma $h(x)\exp(\eta T(x)-A(\eta))$; identifique $\eta$, $T$, $A$ y $h$. \textit{Pista:} la densidad es $\beta^\alpha/\Gamma(\alpha)\,x^{\alpha-1}e^{-\beta x}=\exp((\alpha{-}1)\log x\cdot (\cdots))$.
	\item[$\star\star\star$] \textbf{PKD para Bernoulli.} Demuestre por qué basta $T(x)=x$ como estadística suficiente para $\mathrm{Bern}(p)$. \textit{Hint:} factorice la verosimilitud de $n$ muestras en $g(T(x),p)\cdot h(x)$.
	\item[$\star\star\star$] \textbf{Softmax es exponencial.} Muestre que la categórica $K$-aria $\mathrm{Categorical}(\pi_1,\dots,\pi_K)$ con $\pi_i=e^{\eta_i}/\sum_j e^{\eta_j}$ es familia exponencial con $\eta=(\eta_1,\dots,\eta_{K-1})^T$, $T(x)=(e_{x_1},\dots,e_{x_{K-1}})^T$, $A(\eta)=\log\sum_j e^{\eta_j}$.
	\item[$\star\star\star\star$] \textbf{MLE vs. moment matching.} Implemente máxima verosimilitud y moment-matching en mezclas de 2 normales univariadas con 1000 muestras sintéticas y compare sesgo/viés cuadrático medio en cada componente. \textit{Discuta:} ¿cuándo difieren, y por qué PKD no aplica?
	\item[$\star\star\star$] Muestre que la familia de Cauchy $p_\theta(x)=1/(\pi(1+(x-\theta)^2))$ \emph{no} es familia exponencial (sin estadístico suficiente de dimensión finita). \textit{Hint:} $\log p_\theta=-\log\pi-\log(1+(x-\theta)^2)$ no es lineal en $T(x)$.
\end{enumerate}

\subsection*{Ejercicios avanzados}
\begin{enumerate}[resume]
	\item[$\star$] (V/F) La $\alpha$-conexión tiene torsión si y solo si la familia no es dualmente plana. \textit{Pregunta a tu LLM:} «¿Por qué la dual planitud hace que $\Gamma^{(1)}=\Gamma^{(-1)}$?»
	\item[$\star\star$] En $\mathrm{Bern}(p)$, grafique las tres $\alpha$-geodésicas ($\alpha=-1,0,1$) entre $p=0.2$ y $p=0.8$. \textit{Verifica con:} la $\alpha=1$ pasa por $\eta=(\log 0.2/0.8+\log 0.8/0.2)/2=0$ (punto medio $p=0.5$).
	\item[$\star\star\star$] Demuestre que $T^{(\alpha)}=\alpha(\Gamma^{(1)}-\Gamma^{(-1)})$. \textit{Pista:} escriba los símbolos de Christoffel explícitamente y observe el cambio al permutar índices $j,k$.
	\item[$\star\star\star$] Demuestre que las $\alpha$-geodésicas en una familia exponencial satisfacen $\ddot\eta+\Gamma^{(\alpha)}\dot\eta^2=0$ (Riccati). \textit{Hint:} aplique $\alpha$-conexión a la geodésica en coordenadas $\eta$.
	\item[$\star\star\star\star$] \textbf{Mini-proy.} Implemente el cálculo numérico de geodésicas $\alpha$ en $(\mu,\sigma)$ entre dos Gaussianas (RK4 sobre la ODE geodésica). Represente las tres variantes en 2D. \textit{Discuta:} ¿cuál cruza el ``infinito'' $\sigma=0$?
\end{enumerate}

\subsection*{Ejercicios avanzados}
\begin{enumerate}[resume]
	\item[$\star$] (V/F) Toda distribución tiene un estadístico suficiente de dimensión finita. \textit{Pregunta a tu LLM:} «¿Por qué Cauchy rompe esa intuición?»
	\item[$\star\star$] Calcule $I(\lambda)$ para $\mathrm{Pareto}(\theta)$ con $\theta$ desconocido y densidad $\theta x^{-\theta-1}$. \textit{Verifica con:} $I(\theta)=1/\theta^2$.
	\item[$\star\star\star$] Para $\mathrm{Poisson}(\lambda)$ con $n=50$ muestras, derive $\operatorname{Var}(\hat\lambda)$ para MLE y compárelo con la cota $1/(50 I(50))$. \textit{Verifica con:} similación con $10^4$ repeticiones.
	\item[$\star\star\star$] Demuestre que la familia de Cauchy con parámetro de localización $\theta$ no admite estadístico suficiente de dimensión finita. \textit{Hint:} afecte a dos puntos $x,y$: el likelihood depende de ellos de manera acoplada.
	\item[$\star\star\star\star$] \textbf{Mini-proy.} Muestre empíricamente que la métrica ANY covariante invariante bajo Markov morphisms es proporcional a Fisher. Use Bernoulli: genere 10000 muestras, calcule $I(p)$ bajo 3 Markov morphisms diferentes (subsampling, ruido simétrico, binning) y reporte la razón.
\end{enumerate}

\subsection*{Ejercicios avanzados}
\begin{enumerate}[resume]
	\item[$\star$] Defina la entropía cruzada binaria en términos de Bregman. \textit{Pregunta a tu LLM:} «¿Qué $F$ da $H(p,q)=-p\log q-(1-p)\log(1-q)$?»
	\item[$\star\star$] Para Bühlmann con EPV$=2$ y VHM$=8$ calcule $Z$ para $n=20$. \textit{Verifica con:} $Z=20/(20+\text{EPV}/\text{VHM})=20/24\approx 0.833$.
	\item[$\star\star\star$] Muestre que la pérdida de entropía cruzada binaria es una divergencia de Bregman con $F(x)=-x\log x-(1-x)\log(1-x)$. \textit{Pista:} verifique $D_F(p\|q)=H(p,q)$.
	\item[$\star\star\star$] Para VI con prior $\mathcal{N}(0,\tau^2)$ y likelihood $\mathcal{N}(\theta,1)$, derive $\mu_q,\sigma_q$ para $n$ datos. \textit{Verifica con:} posterior exacto igual por conjugación.
	\item[$\star\star\star\star$] \textbf{Mini-proy GLM (estrella del libro).} Implemente regresión Poisson con offset, regularization L2; compare: GD sobre $\eta$, GD sobre $\mu$, Newton sobre ambos, gradiente natural. Reporte convergencia (iteraciones a tol=1e-8) sobre un dataset sintético ($n=1000$). \textit{Discuta:} ¿cuál converge más rápido?
\end{enumerate}

\subsection*{Ejercicios avanzados}
\begin{enumerate}[resume]
	\item[$\star\star$] Calcule la métrica de Fisher de $\mathrm{Cauchy}$ con localización $\theta$. \textit{Pista:} use $\partial_\theta \ell = 2(x-\theta)/(1+(x-\theta)^2)$ y tome esperanza bajo $P_\theta$.
	\item[$\star\star\star$] Para un qubit $|\psi\rangle=\cos(\theta/2)|0\rangle+\sin(\theta/2)e^{i\phi}|1\rangle$, la métrica de Bures induce $ds^2=(d\theta^2+\sin^2\theta\,d\phi^2)/4$; identifique su tipo métrico ($S^2$?). \textit{Verifica con:} $\int\theta\in[0,\pi],\phi\in[0,2\pi]$ cubre la esfera.
	\item[$\star\star\star$] \textbf{Wasserstein 1D.} Para distribuciones discretas $p,q$ con soporte $\{x_1<\cdots<x_k\}$, demuestre $W_1(p,q)=\sum_{i=1}^{k-1}|F_p(x_i)-F_q(x_i)|(x_{i+1}-x_i)$, donde $F$ es la FDA. \textit{Pista:} el acoplamiento óptimo ordena ambas FDA y aplica ``northwest corner''.
	\item[$\star\star\star$] \textbf{T-Student =?} Demuestre que la familia t-Student con $\nu$ g.l.\\ y localización $\mu$ no tiene estadístico suficiente de dimensión finita. \textit{Verifica con:} la verosimilitud conjunta de $n$ muestras no factoriza.
	\item[$\star\star\star\star$] \textbf{Mini-proy Wasserstein.} Entrene un flow normalizado simple (NICE: 4 capas de acoplamiento affine + permutaciones) sobre MNIST $8\times 8$ con pérdida Wasserstein-2. Reporte log-likelihood aproximado y muestras visualmente. \textit{Discuta:} ¿cómo el costo Wasserstein cambia la dinámica de entrenamiento?
\end{enumerate}

\subsection*{Qué gana con GI}
\begin{quote}
	\itshape
	Sin GI, la métrica de Fisher se define como la varianza del score y se usa como herramienta técnica.  Con GI, se entiende como la \emph{única} métrica invariante bajo transformaciones que preservan información (Chentsov).  Esto significa que cualquier conclusión derivada de Fisher--Rao---como la cota de Cramér--Rao o el prior de Jeffreys---es \emph{independiente de cómo parametrices el modelo}.  El resultado no es cuestión de gusto: es una propiedad geométrica del espacio de distribuciones.
\end{quote}
